% ============================================================================
% Chapter 5 — Modern Arabic Poetry Forms
% ============================================================================
\chapter{Modern Arabic Poetry Forms}
\label{ch:modern}

The mid-twentieth century brought a revolution to Arabic poetry that severed — or at least dramatically loosened — the bond between poetry and the classical prosodic system. Two new forms emerged: \textit{shi'r al-tafila} (\ar{شعر التفعيلة}), which retains classical meter but abandons the fixed two-hemistich line; and \textit{qasidat al-nathr} (\ar{قصيدة النثر}), which abandons both meter and rhyme in favour of concentrated prose-rhythm. Understanding these forms is essential for \poeteng{} because they are widely requested and their evaluation criteria differ fundamentally from classical and Nabati poetry.

% ============================================================================
\section{\ar{شعر التفعيلة} — Metrical Free Verse}
\label{sec:tafila}

\subsection{Historical Origins}

The emergence of \ar{شعر التفعيلة} is dated to December 1947, when the Iraqi poet \textbf{Nazik al-Mala'ika} (\ar{نازك الملائكة}, 1923–2007) published her poem ``Cholera'' (\ar{الكوليرا}) in the Lebanese journal \textit{al-Uruba}. The poem used classical metrical feet (\ar{تفعيلات}) but in lines of variable length — a radical departure from the fixed-line bayt.

The Iraqi poet \textbf{Badr Shakir al-Sayyab} (\ar{بدر شاكر السياب}, 1926–1964) published ``Was It About Love?'' at nearly the same time and is co-credited as founder. Other early pioneers: \textbf{Abd al-Wahhab al-Bayati} and \textbf{Shathel Taqa}. Nazik al-Mala'ika later theorised the form in her landmark study \textit{Issues of Modern Poetry} (\ar{قضايا الشعر المعاصر}, 1962).

Major later poets: Fadwa Tuqan, Adonis, Mahmoud Darwish, Amal Dunqul, Ahmad Abd al-Mu'ti Hijazi, Samih al-Qasim.

\subsection{Structural Rules}

\begin{rulebox}{Defining Characteristics of \ar{شعر التفعيلة}}
\begin{enumerate}
  \item \textbf{Retains classical quantitative meter}: The poem is built on one chosen Khalilian bahr's \textit{tafila} (foot). Valid choices: Kamil, Mutaqarib, Ramal, Wafir, Hazaj, Baseet.
  \item \textbf{Variable line length}: Each line (\ar{سطر}) may contain one, two, three, or more iterations of the chosen foot — the number is determined by the poem's meaning and emotion, not by a fixed scheme.
  \item \textbf{Rhyme is optional}: Rhyme may appear but is never enforced. When it appears, it is selective (not monorhyme). End-rhyme is possible but not mandatory.
  \item \textbf{The satra (line) replaces the bayt}: The unit of composition is the line (\ar{السطر}), which can be short (one foot) or long (many feet). Lines are not paired into hemistichs.
  \item \textbf{The tafila is the binding element}: The poem's internal musicality comes from the repeated foot, not from structural regularity.
\end{enumerate}
\end{rulebox}

\subsection{Meter Selection in \ar{شعر التفعيلة}}

Not all meters work equally well in تفعيلة poetry. Nazik al-Mala'ika and other theorists observed:

\begin{table}[H]
\centering
\caption{Meter suitability for \ar{شعر التفعيلة}}
\renewcommand{\arraystretch}{1.5}
\begin{tabularx}{\textwidth}{llX}
\toprule
\textbf{Meter} & \textbf{Suitability} & \textbf{Notes} \\
\midrule
\ar{الكامل} (Kamel) & Excellent & The most emotionally versatile; used extensively by Darwish \\
\ar{المتقارب} (Mutakareb) & Excellent & Light, flowing; used by al-Sayyab in epic poems \\
\ar{الرمل} (Ramal) & Very good & Lyrical, melancholic \\
\ar{الوافر} (Wafir) & Good & Emotional depth \\
\ar{الهزج} (Hazaj) & Good & Musical, light \\
\ar{البسيط} (Baseet) & Moderate & Heavier feel; less common in تفعيلة \\
\ar{الطويل} (Taweel) & Difficult & The long complex feet create monotony at variable lengths \\
\ar{الرجز} (Rajaz) & Special case & Accepted in didactic/experimental تفعيلة \\
\bottomrule
\end{tabularx}
\end{table}

\subsection{Al-Sayyab's ``Rain Song'' — A Reference Example}

\textbf{Badr Shakir al-Sayyab's ``Unshoudat al-Matar'' (\ar{أنشودة المطر})} is the canonical example of تفعيلة poetry. It uses the \textbf{Kamel meter} (\ar{مُتَفَاعِلُنْ}) with variable line lengths:

\begin{verseframe}
\centering
\ar{عَيْنَاكِ غَابَتَا نَخِيلٍ سَاعَةَ السَّحَرْ} \\[4pt]
\ar{أَوْ شُرْفَتَانِ رَاحَ يَنْأَى عَنْهُمَا الْقَمَرْ} \\[4pt]
\ar{عَيْنَاكِ حِينَ تَبْتَسِمَانِ تُورِقُ الْكُرُومْ} \\[4pt]
\ar{وَتَرْقُصُ الأَضْوَاءُ كَالأَقْمَارِ فِي نَهَرْ} \\[4pt]
\ar{يَهُزُّهُ الْمِجْدَافُ وَهْنًا سَاعَةَ السَّحَرْ}\\[4pt]
\small\textit{Al-Sayyab, ``Rain Song'' (opening) — Kamel meter, variable line lengths}
\end{verseframe}

\subsection{Evaluation Criteria for \ar{شعر التفعيلة}}

The evaluator must check:
\begin{enumerate}
  \item \textbf{Metrical consistency}: Every line must be constructible from valid iterations of the chosen foot, allowing standard zehafs.
  \item \textbf{Rhyme logic} (if present): If the poet has introduced rhyme, it should be consistent within its pattern — not random.
  \item \textbf{Internal rhythm}: Even without fixed line length, the poem should have a sense of rhythmic flow — varying line lengths deliberately for effect.
  \item \textbf{Imagery quality}: The standard for modern poetry is judged heavily by imagery, metaphor, and emotional resonance.
  \item \textbf{Lineation logic}: Line breaks should be meaningful — not arbitrary. A line should break where a musical or semantic beat falls.
\end{enumerate}

\begin{implnote}
For AI validation of \ar{شعر التفعيلة}: the classical prosody library (bohour) can be applied at the level of individual lines rather than bayt pairs. Each line must scan as a valid sequence of the chosen tafila with allowed zehafs. Line pairs do not need to match each other's length.
\end{implnote}

% ============================================================================
\section{\ar{قصيدة النثر} — The Prose Poem}
\label{sec:prose_poem}

\subsection{Historical Origins and Theoretical Grounding}

\ar{قصيدة النثر} completely abandons both meter and rhyme. It was theoretically launched in 1960 through two simultaneous landmark texts:

\begin{itemize}
  \item \textbf{Adonis} (\ar{أدونيس}, Ali Ahmad Sa'id Esber): published ``On the Prose Poem'' (\ar{في قصيدة النثر}), proclaiming it ``a distinctive artistic construction with its own identity.''
  \item \textbf{Unsi al-Hajj} (\ar{أنسي الحاج}): published his collection \textit{Lan} (\ar{لن}) with a seminal theoretical introduction — the first use of \ar{قصيدة النثر} as a formal category.
\end{itemize}

Both texts drew theoretical inspiration from Suzanne Bernard's 1959 French study of the prose poem. The institutional home was \textbf{Majallat Shi'r} (\ar{مجلة شعر} — ``Poetry Magazine''), founded in Beirut in 1957 by Yusuf al-Khal (\ar{يوسف الخال}).

\textbf{Nazik al-Mala'ika} — who championed تفعيلة poetry — actively \textit{opposed} the prose poem, viewing it as either a misnomer or incompetent verse. This opposition is itself important context: the prose poem poses a radical challenge to the definition of Arabic poetry as ``metered and rhymed speech'' (\ar{كلام موزون مقفى}).

\subsection{Defining Characteristics}

\begin{rulebox}{What Makes a \ar{قصيدة نثر}}
\begin{enumerate}
  \item \textbf{No quantitative meter}: The text cannot be divided into classical feet.
  \item \textbf{No rhyme}: End-rhyme is absent (though internal assonance may exist).
  \item \textbf{Concentrated language}: Dense imagery, high semantic load per word.
  \item \textbf{Internal rhythm}: The ``music'' comes from rhythm of syntax, repetition patterns, vowel/consonant textures — not prosodic pattern.
  \item \textbf{Unity}: A prose poem is shorter than conventional prose; it has a single emotional or conceptual arc.
  \item \textbf{Distinctness from prose}: The text is richer in imagery and denser in language than ordinary literary prose.
\end{enumerate}
\end{rulebox}

\subsection{Evaluation Criteria for \ar{قصيدة نثر}}

Evaluation of prose poetry cannot rely on prosodic validation. Criteria:

\begin{table}[H]
\centering
\caption{Prose poem evaluation dimensions}
\renewcommand{\arraystretch}{1.5}
\begin{tabularx}{\textwidth}{lX}
\toprule
\textbf{Dimension} & \textbf{Criteria} \\
\midrule
Imagery density & Are the images fresh and unexpected? Is there metaphoric depth? \\
Language concentration & Every word should be load-bearing; padding/fillers are defects. \\
Internal rhythm & Syntactic patterns, deliberate repetition, and phonological texture create rhythm. \\
Emotional arc & Does the poem have a beginning, development, and resolution — even if non-narrative? \\
Distinctness from prose & Is this genuinely ``poetic'' — or just paragraph writing? \\
Unity & Is there a single governing image or concept that holds the piece together? \\
\bottomrule
\end{tabularx}
\end{table}

\begin{implnote}
For AI generation of \ar{قصيدة نثر}: the system should use a different prompt paradigm than for classical/Nabati poetry. No tafeelat, no rhyme constraints. The generation goal is: concentrated imagery + emotional arc + language density. The evaluator cannot use the bohour library — it should use LLM-based quality assessment only.
\end{implnote}

% ============================================================================
\section{Mixing Traditions (\ar{التمزيج بين التقاليد})}
\label{sec:mixing}

Modern Arabic poets frequently work at the intersection of traditions. Some patterns:

\begin{itemize}
  \item \textbf{Neo-classical (\ar{الإحيائي})}: Classical meters and forms but modern content. Ahmad Shawqi (``the Prince of Poets'') and Hafiz Ibrahim. This is technically still ``classical'' in the system's taxonomy.
  \item \textbf{Romantic (\ar{الرومانسي})}: Influenced by Western Romanticism; uses classical forms but with freer imagery and more personal emotional content. Khalil Gibran's Arabic prose poetry; the Mahjar school.
  \item \textbf{Engaged poetry (\ar{الشعر الملتزم})}: Political and social commitment fused with modern forms. Mahmoud Darwish (Palestine), Nizar Qabbani (Syria — who used classical meters for provocative modern content), Amal Dunqul.
  \item \textbf{Hybrid classical/Nabati}: Some contemporary poets compose in both traditions, or use Nabati meters with Fusha language. This is technically non-standard but requested.
\end{itemize}

\begin{table}[H]
\centering
\caption{Modern Arabic poetry tradition overview for AI}
\renewcommand{\arraystretch}{1.5}
\begin{tabularx}{\textwidth}{lllX}
\toprule
\textbf{Form} & \textbf{Meter?} & \textbf{Rhyme?} & \textbf{AI Evaluation Approach} \\
\midrule
Classical qasida & Yes (Khalilian) & Yes (monorhyme) & Bohour library + qafiya rules \\
Neo-classical & Yes (Khalilian) & Yes & Same as classical \\
\ar{شعر التفعيلة} & Yes (single foot) & Optional & Bohour library per-line \\
\ar{قصيدة النثر} & No & No & LLM qualitative only \\
Nabati (Mashhub) & Yes (Nabati) & Yes (dual) & Nabati-specific rules \\
Zajal & Yes (stress-syllabic) & Yes (strophic) & Form-specific rules \\
\bottomrule
\end{tabularx}
\end{table}

% ============================================================================
\clearpage
\langsep{}

% ============================================================================
%  ARABIC VERSION
% ============================================================================
\begin{otherlanguage}{arabic}

\section*{الأشكال الشعرية الحديثة}

\subsection*{شعر التفعيلة}

ظهر شعر التفعيلة في العراق أواخر الأربعينيات من القرن العشرين. تُنسَب ريادته إلى نازك الملائكة (1923–2007) التي نشرت قصيدة ``الكوليرا'' في ديسمبر 1947، وإلى بدر شاكر السياب (1926–1964) الذي نشر قصيدة ``هل كان حبًّا؟'' في نفس الفترة تقريبًا.

\textbf{مبادئ شعر التفعيلة:}
\begin{enumerate}
  \item يحتفظ بالبحور الخليلية (تفعيلاتها)، لكن يتخلى عن الشطرَين الثابتَي الطول.
  \item الوحدة هي السطر (\ar{السطر}) لا البيت؛ وطول السطر متغير حسب المعنى والإيقاع.
  \item القافية اختيارية وليست ملزمة؛ إذا وُجدت فهي انتقائية لا متواترة.
  \item التفعيلة المختارة هي العمود الفقري الموسيقي للقصيدة.
\end{enumerate}

\textbf{أفضل البحور لشعر التفعيلة:} الكامل (محمود درويش) والمتقارب (السياب في قصائده الملحمية) والرمل والوافر.

\textbf{قواعد تقييم الذكاء الاصطناعي:} يتحقق النظام من أن كل سطر يتكون من تكرارات صحيحة للتفعيلة المختارة بزحافاتها المسموحة — لكنه لا يشترط تساوي أطوال الأسطر.

\textbf{مثال} (من ``أنشودة المطر'' للسياب على بحر الكامل):
عَيْنَاكِ غَابَتَا نَخِيلٍ سَاعَةَ السَّحَرْ --- أَوْ شُرْفَتَانِ رَاحَ يَنْأَى عَنْهُمَا الْقَمَرْ

\subsection*{قصيدة النثر}

قصيدة النثر تتخلى كليًّا عن الوزن والقافية. أطلقها نظريًا أدونيس وأنسي الحاج عام 1960 في مجلة ``شعر'' البيروتية.

\textbf{ما يميز قصيدة النثر عن النثر العادي:}
\begin{itemize}
  \item كثافة الصور الشعرية وثقل الدلالة في كل كلمة.
  \item إيقاع داخلي ناتج عن بنية الجملة والتكرار والتوازي اللفظي.
  \item وحدة موضوعية أو صورية تجمع القصيدة.
  \item إيجاز لا يُتاح للنثر الاعتيادي.
\end{itemize}

\textbf{قواعد تقييم الذكاء الاصطناعي:} لا يُطبَّق النظام العروضي؛ يُستخدم تقييم نوعي مبني على الذكاء الاصطناعي يرصد: كثافة الصور، عمق المجاز، الإيقاع الداخلي، القوس العاطفي، والكثافة اللغوية.

\subsection*{المزج بين التقاليد}

بعض الأنماط الشعرية الحديثة تجمع بين التقاليد:
\begin{description}
  \item[\textbf{الإحيائي}] شعراء النهضة كأحمد شوقي وحافظ إبراهيم: مضمون حديث في قوالب كلاسيكية.
  \item[\textbf{الرومانسي}] مدرسة المهجر (جبران خليل جبران): نثر شعري يمزج العربية بالرومانسية الغربية.
  \item[\textbf{الملتزم}] محمود درويش ونزار قباني وأمل دنقل: التزام سياسي واجتماعي في أشكال حديثة.
\end{description}

\end{otherlanguage}
